\documentclass[11pt]{article}
\usepackage[a4paper,margin=1in]{geometry}
\usepackage[T1]{fontenc}
\usepackage[utf8]{inputenc}
\usepackage{lmodern}
\usepackage{microtype}
\usepackage{amsmath,amssymb,amsfonts,bm}
\usepackage{graphicx}
\usepackage{booktabs}
\usepackage{float}
\usepackage{caption}
\usepackage{setspace}
\usepackage{hyperref}
\usepackage{enumitem}
\usepackage{mathtools}
\usepackage{cite} % 引用番号を[1-3]の形式にまとめるために追加

\hypersetup{colorlinks=true, linkcolor=black, citecolor=black, urlcolor=blue}
\setstretch{1.08}
\setlength{\parskip}{0.45em}
\setlength{\parindent}{1.2em}
\captionsetup{font=small,labelfont=bf}

\title{\vspace{-1.5em}\textbf{Socio-Traffic Thermodynamics:}\\[0.25em]
\large Congestion as a Non-Equilibrium Dissipative Structure Driven by Social Potential}
\author{Keiji Yoshimura \\ \small Independent Researcher}
\date{April 16, 2026\\[0.35em]
\small Repository DOI: \href{https://doi.org/10.5281/zenodo.20201218}{10.5281/zenodo.20201218}\\
\small GitHub repository: \href{https://github.com/yokken0907/socio-traffic-thermodynamics}{github.com/yokken0907/socio-traffic-thermodynamics}}

\begin{document}
\maketitle

\begin{abstract}
We propose Socio-Traffic Thermodynamics (STT), a theoretical framework that models traffic congestion as a non-equilibrium dissipative structure generated by agents moving along social-potential gradients on transportation networks with finite capacity, finite perceptual bandwidth, and finite relaxation time. The framework integrates three mechanisms: bandwidth compression induced by vehicular encapsulation, demand refilling under endogenous cost equilibration, and attraction toward high-opportunity, high-interaction regions. We formulate a microscopic model combining stochastic car-following dynamics with an internal binary state field, and derive a macroscopic continuum model with density, mean velocity, and an order parameter representing synchronized internal orientation. We show that near the free-flow critical point, decentralized residual control variance---interpreted as sovereignty fluctuation---monotonically increases dissipation and lowers the effective barrier against jam nucleation. We further show that coarse or privatized information can improve system performance by reducing action correlation and load variance, thereby lowering expected social cost under convex congestion functions. Numerical simulation on a periodic ring using a stochastic Optimal Velocity model exhibits spontaneous symmetry breaking and persistent backward-propagating stop-and-go waves without any exogenous bottleneck. STT does not claim that congestion is universally eliminable. Instead, it characterizes the structural recurrence of congestion in socially attractive, capacity-limited systems and suggests a transition from control-centered traffic engineering to buffering-centered architecture. Within this view, congestion is the irreversible cost paid by a socially driven mobile population when traversing a finite-bandwidth network.
\end{abstract}

\section{Introduction}
Traffic congestion has traditionally been treated as an engineering inefficiency to be minimized by infrastructure expansion, signal control, or demand management. Two robust findings complicate this framing. First, congestion can arise endogenously in the absence of an explicit bottleneck: small fluctuations in a homogeneous stream grow, destabilize free flow, and generate backward-propagating jam waves once density exceeds a critical threshold \cite{Bando1995, Sugiyama2008, Helbing2001}. Second, traffic instabilities are amplified by stochasticity and by the loss of string stability, indicating that congestion is not reducible to static capacity shortage alone.

At the network level, selfish routing generally fails to minimize aggregate latency; the gap is formalized by the price of anarchy (PoA) \cite{Roughgarden2003}. Additional route information is not uniformly beneficial: informational Braess-type effects can make better-informed users worse off on some networks \cite{Acemoglu2018}, while carefully designed private signals can, in some classes of routing games, move equilibrium flows toward the social optimum \cite{Das2017}. These observations motivate a different question: not how to describe congestion as a recurrent macroscopic pattern produced by socially driven motion on finite-bandwidth networks.

STT addresses this question by treating congestion as a non-equilibrium dissipative structure rather than a mere operational failure. The aim is not to replace traffic theory with metaphor, but to provide a mesoscopic language that connects social attraction, microscopic fluctuation, and macroscopic breakdown within one formalism.

\section{Model Formulation}
\subsection{Network and state variables}
Let $G=(V,E)$ be a directed transportation network. Let $\Omega$ denote a continuum or a large finite population of agents. Each agent $i\in\Omega$ is characterized by position $x_i(t)$, velocity $v_i(t)$, internal state $s_i(t)\in[-1,1]$, route $p_i(t)\in\mathcal P$, and destination-dependent social potential $\Phi_i = \Phi(x_i,t)$.

Define the macroscopic fields
\begin{equation}
\rho(x,t),\qquad u(x,t),\qquad m(x,t):=\mathbb E[s_i\mid x_i=x],
\end{equation}
where $\rho$ is density, $u$ is mean velocity, and $m$ is an order parameter encoding local synchronization of internal states.

\subsection{Microscopic longitudinal dynamics}
For each agent $i$,
\begin{equation}
\dot x_i=v_i,
\end{equation}
\begin{equation}
\dot v_i = F(h_i,\Delta v_i,s_i;\rho_i,\Phi_i) + \varepsilon_i(t) + \sqrt{2D_v}\,\xi_i(t),
\end{equation}
where $h_i$ is the headway, $\Delta v_i$ is the relative velocity, $F$ is a smooth car-following response, $\varepsilon_i(t)$ is decentralized residual control, and $\xi_i(t)$ is a white-noise term.

The internal-state dynamics is given by
\begin{equation}
\tau_s \dot s_i = -s_i + \tanh\!\Bigl(\beta\bigl[J(\mathcal K*s)_i + \lambda \Phi_i - \kappa C_i\bigr]\Bigr) + \sqrt{2D_s^{\mathrm{eff}}}\,\zeta_i(t),
\end{equation}
with perceived congestion cost $C_i$ and effective internal noise
\begin{equation}
D_s^{\mathrm{eff}} = D_{s,0} + \chi_{\mathrm{sov}}\sigma_{\mathrm{sov}}^2 + \chi_W W_i^{-1}.
\end{equation}
Here $W_i$ denotes effective perceptual bandwidth. Equation (5) formalizes the claim that decentralized idiosyncrasy and information compression both act as effective temperature.

\subsection{Route choice}
Given information $\mathcal I_i$, agent $i$ chooses a path $p\in\mathcal P$ with probability
\begin{equation}
\mathbb P(p_i=p\mid \mathcal I_i)=
\frac{\exp\{-\theta[\mathcal C_p-\lambda_\Phi \Phi_p]\}}
{\sum_{q\in\mathcal P}\exp\{-\theta[\mathcal C_q-\lambda_\Phi \Phi_q]\}},
\end{equation}
where $\mathcal C_p$ is generalized travel cost along path $p$ and $\Phi_p$ is destination-adjusted social potential.

\subsection{Macroscopic closure}
Under mean-field closure, the macroscopic dynamics is
\begin{equation}
\partial_t \rho + \nabla\!\cdot (\rho u)=0,
\end{equation}
\begin{equation}
\partial_t(\rho u)+\nabla\!\cdot(\rho u\otimes u)= -\nabla P(\rho,m)-\frac{\rho}{\tau}\bigl(u-U(\rho,m,\Phi)\bigr)+\nu\Delta u + \sqrt{2D_u\rho}\,\eta,
\end{equation}
\begin{equation}
\tau_m\partial_t m = A(\rho-\rho_c)m - Bm^3 + D_m\Delta m + \lambda\Phi - \kappa C(\rho) - m + \sqrt{2D_m^{\mathrm{eff}}}\,\zeta.
\end{equation}

\section{Metastability and Critical Transition}
\subsection{Linearization and metastability}
Let $(\rho_0,u_0,m_0)$ be a homogeneous free-flow equilibrium. Fourier linearization yields
\begin{equation}
\dot{\hat z}_k = A_k\hat z_k + B_k\hat \omega_k + G_k\hat \xi_k.
\end{equation}
The system is linearly stable if $\Re\lambda(A_k)<0$ for all $k$. It is \emph{metastable} if this remains true while the stability margin
\begin{equation}
\alpha := \min_k\{-\Re\lambda_{\max}(A_k)\}
\end{equation}
is small and positive. In this regime, perturbations decay only slowly, susceptibility is large, and rare transitions become likely.

\subsection{Cost of sovereignty}
Let $P_k$ be the stationary covariance of mode $k$, satisfying
\begin{equation}
A_kP_k + P_kA_k^{\ast} + Q_k = 0,
\qquad
Q_k = Q_{0,k} + \sigma_{\mathrm{sov}}^2 B_kB_k^{\ast}.
\end{equation}
For any positive semidefinite weight $R_k$, define
\begin{equation}
\Pi_k = \mathrm{tr}(R_kP_k).
\end{equation}
Then
\begin{equation}
\frac{\partial \Pi_k}{\partial \sigma_{\mathrm{sov}}^2}
=
\mathrm{tr}\!\left(R_k\mathcal L_k^{-1}(B_kB_k^{\ast})\right) \ge 0.
\end{equation}
Hence total dissipation $\Pi = \sum_k \Pi_k$ is nondecreasing in $\sigma_{\mathrm{sov}}^2$. Moreover,
\begin{equation}
\Pi \gtrsim \frac{c\,\sigma_{\mathrm{sov}}^2}{\alpha}, \qquad \alpha\to 0^+.
\end{equation}
This states that decentralized residual freedom becomes macroscopically expensive precisely where the system loses stability margin.

\subsection{Nucleation of congestion}
Let $a$ denote the amplitude of the critical stop-and-go mode. Near the transition,
\begin{equation}
\dot a = -\partial_a \mathcal V(a;\rho,\Phi) + \sqrt{2D_{\mathrm{eff}}}\,\eta(t),
\end{equation}
where $\mathcal V$ is an effective quasi-potential and
\begin{equation}
D_{\mathrm{eff}} = D_0 + \chi_{\mathrm{sov}}\sigma_{\mathrm{sov}}^2 + \chi_W W^{-1}.
\end{equation}
The nucleation rate of a jam wave is asymptotically
\begin{equation}
r_{\mathrm{nuc}} \asymp \exp\!\left(-\frac{\Delta \mathcal V(\rho,\Phi)}{D_{\mathrm{eff}}}\right).
\end{equation}

\section{Information Buffering}
\subsection{Signal structure}
Let $X$ denote network state and $Y_i$ the signal received by agent $i$. A signal policy is a stochastic kernel $K(y\mid x)$. Path choices induced by $\{Y_i\}$ generate edge loads
\begin{equation}
F_e = \sum_{i=1}^{N} \mathbf 1\{e\in A_i\}.
\end{equation}
The expected social cost is
\begin{equation}
\mathbb E[SC(F)] = \sum_{e\in E} \mathbb E[F_e c_e(F_e)].
\end{equation}

\subsection{Variance principle}
If $SC_e(f)=fc_e(f)$ is convex, then
\begin{equation}
\mathbb E[SC(F)] = SC(\bar F) + \frac{1}{2}\sum_e SC_e''(\bar F_e)\,\mathrm{Var}(F_e) + o\bigl(\mathrm{Var}(F_e)\bigr).
\end{equation}
Further,
\begin{equation}
\mathrm{Var}(F_e)=\sum_i \mathrm{Var}(A_{ie}) + 2\sum_{i<j}\mathrm{Cov}(A_{ie},A_{je}).
\end{equation}

\subsection{Utility of ignorance}
A public high-resolution signal tends to reduce $H(X\mid Y)$, but it also tends to increase conditional action correlation $I(A_i;A_j\mid X)$, which raises the covariance term in Eq.~(19). By contrast, garbled or privatized signals can preserve enough state adaptation while reducing action synchronization.

\subsection{PoA under information design}
For a signaling rule $K$, define
\begin{equation}
\mathrm{PoA}(K)=\frac{\mathbb E[SC(F^{K}_{\mathrm{eq}})]}{\mathbb E[SC(F^{\star})]}.
\end{equation}
Information design acts primarily on the numerator by reshaping the equilibrium covariance structure. Additional information can worsen equilibrium outcomes on some networks \cite{Acemoglu2018}, while private or carefully garbled signals can improve them \cite{Das2017}.

\section{Numerical Validation}
\subsection{Simulation setting}
To examine whether the STT transition mechanism is realizable in a minimal microscopic system, we consider a one-lane periodic ring using an Optimal Velocity (OV) car-following law \cite{Bando1995} with additive Gaussian decentralized residual input. The simulation follows
\begin{equation}
\dot x_i=v_i,
\end{equation}
\begin{equation}
\dot v_i = a\bigl[V(h_i)-v_i\bigr] + \sigma_{\mathrm{sov}}\xi_i(t),
\end{equation}
where $x_i$ is position, $v_i$ is velocity, $h_i=x_{i+1}-x_i$ is headway on the ring, $V(\cdot)$ is the optimal-velocity function, and $\xi_i(t)$ is a zero-mean unit Gaussian process.

\subsection{Space-time observation}
Figure~\ref{fig:spacetime} displays the resulting space-time diagram. At early times, the trajectories are nearly parallel, indicating approximate uniform motion. As the evolution proceeds, localized bunching emerges, the local slope of several trajectories sharply decreases, and a coherent high-density band forms. Once formed, this band persists over long times and propagates backward relative to the direction of vehicle motion. The resulting pattern is the canonical signature of a self-sustained stop-and-go wave \cite{Sugiyama2008}.

\begin{figure}[H]
    \centering
    \includegraphics[width=0.96\textwidth]{stt_result.png}
    \caption{Space-time diagram of the stochastic OV ring-road simulation with periodic boundary conditions and additive sovereignty fluctuation $\sigma_{\mathrm{sov}}$. Each trajectory corresponds to one vehicle. An initially near-homogeneous free-flow state loses symmetry, condenses into a localized high-density cluster, and develops a persistent backward-propagating stop-and-go wave in the absence of any fixed bottleneck.}
    \label{fig:spacetime}
\end{figure}

\subsection{Dissipative-structure interpretation}
The backward-propagating wave should not be interpreted as a conservative traveling mode. It is more appropriately interpreted as a dissipative structure sustained by throughput, relaxation delay, and irreversible adaptation. One convenient operational choice is
\begin{equation}
\Sigma_{\mathrm{irr}} = \int_{\Omega_x} \left[\alpha_1 |\partial_x u|^2 + \alpha_2 |\partial_t u|^2 + \alpha_3 \operatorname{Var}(u\mid x)\right]dx,
\end{equation}
with $u(x,t)$ the coarse-grained velocity field.

Within STT, this is the precise meaning of treating congestion as a dissipative structure. The stop-and-go wave is an organized nonequilibrium pattern whose existence depends on continuous throughput and continuous dissipation \cite{Helbing2001}.

\section{Architectural Implications}
STT implies that the relevant engineering objective is not universal decongestion, but reduction of critical amplification.

\subsection{Temporal buffering}
Design interventions should reduce synchronization gradients $\partial_t\rho$, not merely average demand.

\subsection{Spatial buffering}
Intermediate reservoirs, transfer basins, and pre-network storage reduce abrupt load injection into already metastable regions.

\subsection{Informational buffering}
Routing guidance should be designed as a variance-control channel, not as a universal shortest-path oracle.

\subsection{Behavioral buffering}
A small fraction of damping-capable agents can reduce oscillatory amplification by increasing local relaxation.

\section{Discussion}
STT concerns a driven, dissipative, stochastic, open system. The thermodynamic vocabulary is retained only where an operational correspondence is specified. The ring-road simulation strengthens this claim by showing that even without exogenous defects, a near-uniform stream can collapse into a stable stop-and-go pattern when residual decentralized fluctuation is present near the critical regime \cite{Sugiyama2008}.

\section{Conclusion}
Within STT, traffic congestion is modeled as a non-equilibrium dissipative structure generated when socially motivated agents move across a finite-bandwidth, finite-capacity transport network. The critical transition is governed by the interaction of density, internal-state synchronization, information compression, and decentralized residual control variance.

The numerical experiment added in this study sharpens that claim. In a periodic OV system with no external bottleneck \cite{Bando1995, Sugiyama2008}, additive sovereignty fluctuation is sufficient to destabilize a near-uniform free-flow state and produce a persistent backward-propagating stop-and-go wave. In strictly scientific language: \emph{traffic is the dissipative signature of socially driven motion on a finite network.}


\section*{Data and Code Availability}
The reproducibility package associated with this manuscript is archived on Zenodo with repository DOI \href{https://doi.org/10.5281/zenodo.20201218}{10.5281/zenodo.20201218} and is also available through the GitHub repository \href{https://github.com/yokken0907/socio-traffic-thermodynamics}{github.com/yokken0907/socio-traffic-thermodynamics}. The archived package contains the manuscript PDF and \LaTeX{} source, the original uploaded ring-road simulation script, a cleaned reproducible simulation script, the generated space-time figure, requirements information, claim-boundary documentation, AI-assistance disclosure, and file-manifest materials. The numerical example should be interpreted as a minimal stochastic Optimal Velocity ring-road validation of the STT transition mechanism, not as a calibrated city-scale traffic model.

\section*{Acknowledgments}
The author conceptually designed the study and formulated the core hypotheses. The author acknowledges the use of large language models as computational and theoretical assistants under direct human supervision.

\begin{center}
\footnotesize\texttt{bc1q8xdppt2844w3t2zx9wu8tc05sqym09m64jc2lq}
\end{center}

\begin{thebibliography}{99}

\bibitem{Bando1995}
M. Bando, K. Hasebe, A. Nakayama, A. Shibata, and Y. Sugiyama,
``Dynamical model of traffic congestion and numerical simulation,''
\emph{Physical Review E} \textbf{51}, 1035--1042 (1995).

\bibitem{Sugiyama2008}
Y. Sugiyama, M. Fukui, M. Kikuchi, K. Hasebe, A. Nakayama, K. Nishinari, S.-i. Tadaki, and S. Yukawa,
``Traffic jams without bottlenecks---experimental evidence for the physical mechanism of the formation of a jam,''
\emph{New Journal of Physics} \textbf{10}, 033001 (2008).

\bibitem{Helbing2001}
D. Helbing,
``Traffic and related self-driven many-particle systems,''
\emph{Reviews of Modern Physics} \textbf{73}, 1067--1141 (2001).

\bibitem{Roughgarden2003}
T. Roughgarden,
``The price of anarchy is independent of the network topology,''
\emph{Journal of Computer and System Sciences} \textbf{67}(2), 341--364 (2003).

\bibitem{Acemoglu2018}
D. Acemoglu, A. Makhdoumi, A. Malekian, and A. Ozdaglar,
``Informational Braess' paradox: The effect of information on traffic congestion,''
\emph{Operations Research} \textbf{66}(4), 893--917 (2018).

\bibitem{Das2017}
S. Das, E. Kamenica, and R. Mirka,
``Reducing congestion through information design,''
working paper (2017).


\bibitem{YoshimuraSTTArchive2026}
K. Yoshimura,
\emph{Socio-Traffic Thermodynamics: Repository Archive},
Zenodo (2026).
DOI: \href{https://doi.org/10.5281/zenodo.20201218}{10.5281/zenodo.20201218}.

\end{thebibliography}

\end{document}